% stamps.tex — Reusable TikZ classification stamps and decorative elements
% Loaded by theme .sty files via % stamps.tex — Reusable TikZ classification stamps and decorative elements
% Loaded by theme .sty files via % stamps.tex — Reusable TikZ classification stamps and decorative elements
% Loaded by theme .sty files via % stamps.tex — Reusable TikZ classification stamps and decorative elements
% Loaded by theme .sty files via \input{graphics/stamps}

% ─── Classification stamp (inline) ──────────────────────────────────
% Usage: \classificationstamp{TEXT}
\newcommand{\classificationstamp}[1]{%
  \begin{tikzpicture}[baseline=(text.base)]
    \node[
      draw=scpstamp,
      line width=1.5pt,
      rounded corners=1pt,
      inner sep=3pt,
      font=\sffamily\bfseries\scriptsize\color{scpstamp},
      opacity=0.85
    ] (text) {\MakeUppercase{#1}};
  \end{tikzpicture}%
}

% ─── Rotated classification stamp (for overlays) ────────────────────
% Usage: \rotatedstamp{TEXT}{angle}
\newcommand{\rotatedstamp}[2]{%
  \begin{tikzpicture}
    \node[
      draw=scpstamp,
      line width=2.5pt,
      rounded corners=2pt,
      rotate=#2,
      inner sep=6pt,
      font=\fontsize{18}{20}\selectfont\ttfamily\bfseries\color{scpstamp},
      opacity=0.7
    ] {\MakeUppercase{#1}};
  \end{tikzpicture}%
}

% ─── Document control number ────────────────────────────────────────
% Typewriter-style control number for headers or metadata
% Usage: \controlnumber{DOC-SCP-2026-0001}
\newcommand{\controlnumber}[1]{%
  {\ttfamily\scriptsize\color{scpgray}#1}%
}

% ─── Redaction bar ──────────────────────────────────────────────────
% Black bar for redacted text of arbitrary length
% Usage: \redacted{approximate width, e.g. 3cm}
\newcommand{\redacted}[1]{%
  \tikz[baseline=(bar.base)]{%
    \node[
      fill=scpblack,
      minimum width=#1,
      minimum height=1.1em,
      inner sep=0pt
    ] (bar) {};
  }%
}

% ─── Section divider ────────────────────────────────────────────────
% Red-accented horizontal rule
% Usage: \scpdivider
\newcommand{\scpdivider}{%
  \vspace{6pt}%
  \noindent
  \begin{tikzpicture}
    \draw[scpred, line width=0.8pt] (0,0) -- (\textwidth,0);
    \fill[scpred] (0.5\textwidth-3pt, -2pt) rectangle (0.5\textwidth+3pt, 2pt);
  \end{tikzpicture}%
  \vspace{6pt}%
}

% ─── Clearance level badge ──────────────────────────────────────────
% Usage: \clearancebadge{4}
\newcommand{\clearancebadge}[1]{%
  \begin{tikzpicture}[baseline=(badge.base)]
    \node[
      draw=scpred,
      fill=scpred,
      line width=0.8pt,
      rounded corners=2pt,
      inner xsep=5pt,
      inner ysep=2pt,
      font=\sffamily\bfseries\scriptsize\color{white}
    ] (badge) {LEVEL #1};
  \end{tikzpicture}%
}

\endinput


% ─── Classification stamp (inline) ──────────────────────────────────
% Usage: \classificationstamp{TEXT}
\newcommand{\classificationstamp}[1]{%
  \begin{tikzpicture}[baseline=(text.base)]
    \node[
      draw=scpstamp,
      line width=1.5pt,
      rounded corners=1pt,
      inner sep=3pt,
      font=\sffamily\bfseries\scriptsize\color{scpstamp},
      opacity=0.85
    ] (text) {\MakeUppercase{#1}};
  \end{tikzpicture}%
}

% ─── Rotated classification stamp (for overlays) ────────────────────
% Usage: \rotatedstamp{TEXT}{angle}
\newcommand{\rotatedstamp}[2]{%
  \begin{tikzpicture}
    \node[
      draw=scpstamp,
      line width=2.5pt,
      rounded corners=2pt,
      rotate=#2,
      inner sep=6pt,
      font=\fontsize{18}{20}\selectfont\ttfamily\bfseries\color{scpstamp},
      opacity=0.7
    ] {\MakeUppercase{#1}};
  \end{tikzpicture}%
}

% ─── Document control number ────────────────────────────────────────
% Typewriter-style control number for headers or metadata
% Usage: \controlnumber{DOC-SCP-2026-0001}
\newcommand{\controlnumber}[1]{%
  {\ttfamily\scriptsize\color{scpgray}#1}%
}

% ─── Redaction bar ──────────────────────────────────────────────────
% Black bar for redacted text of arbitrary length
% Usage: \redacted{approximate width, e.g. 3cm}
\newcommand{\redacted}[1]{%
  \tikz[baseline=(bar.base)]{%
    \node[
      fill=scpblack,
      minimum width=#1,
      minimum height=1.1em,
      inner sep=0pt
    ] (bar) {};
  }%
}

% ─── Section divider ────────────────────────────────────────────────
% Red-accented horizontal rule
% Usage: \scpdivider
\newcommand{\scpdivider}{%
  \vspace{6pt}%
  \noindent
  \begin{tikzpicture}
    \draw[scpred, line width=0.8pt] (0,0) -- (\textwidth,0);
    \fill[scpred] (0.5\textwidth-3pt, -2pt) rectangle (0.5\textwidth+3pt, 2pt);
  \end{tikzpicture}%
  \vspace{6pt}%
}

% ─── Clearance level badge ──────────────────────────────────────────
% Usage: \clearancebadge{4}
\newcommand{\clearancebadge}[1]{%
  \begin{tikzpicture}[baseline=(badge.base)]
    \node[
      draw=scpred,
      fill=scpred,
      line width=0.8pt,
      rounded corners=2pt,
      inner xsep=5pt,
      inner ysep=2pt,
      font=\sffamily\bfseries\scriptsize\color{white}
    ] (badge) {LEVEL #1};
  \end{tikzpicture}%
}

\endinput


% ─── Classification stamp (inline) ──────────────────────────────────
% Usage: \classificationstamp{TEXT}
\newcommand{\classificationstamp}[1]{%
  \begin{tikzpicture}[baseline=(text.base)]
    \node[
      draw=scpstamp,
      line width=1.5pt,
      rounded corners=1pt,
      inner sep=3pt,
      font=\sffamily\bfseries\scriptsize\color{scpstamp},
      opacity=0.85
    ] (text) {\MakeUppercase{#1}};
  \end{tikzpicture}%
}

% ─── Rotated classification stamp (for overlays) ────────────────────
% Usage: \rotatedstamp{TEXT}{angle}
\newcommand{\rotatedstamp}[2]{%
  \begin{tikzpicture}
    \node[
      draw=scpstamp,
      line width=2.5pt,
      rounded corners=2pt,
      rotate=#2,
      inner sep=6pt,
      font=\fontsize{18}{20}\selectfont\ttfamily\bfseries\color{scpstamp},
      opacity=0.7
    ] {\MakeUppercase{#1}};
  \end{tikzpicture}%
}

% ─── Document control number ────────────────────────────────────────
% Typewriter-style control number for headers or metadata
% Usage: \controlnumber{DOC-SCP-2026-0001}
\newcommand{\controlnumber}[1]{%
  {\ttfamily\scriptsize\color{scpgray}#1}%
}

% ─── Redaction bar ──────────────────────────────────────────────────
% Black bar for redacted text of arbitrary length
% Usage: \redacted{approximate width, e.g. 3cm}
\newcommand{\redacted}[1]{%
  \tikz[baseline=(bar.base)]{%
    \node[
      fill=scpblack,
      minimum width=#1,
      minimum height=1.1em,
      inner sep=0pt
    ] (bar) {};
  }%
}

% ─── Section divider ────────────────────────────────────────────────
% Red-accented horizontal rule
% Usage: \scpdivider
\newcommand{\scpdivider}{%
  \vspace{6pt}%
  \noindent
  \begin{tikzpicture}
    \draw[scpred, line width=0.8pt] (0,0) -- (\textwidth,0);
    \fill[scpred] (0.5\textwidth-3pt, -2pt) rectangle (0.5\textwidth+3pt, 2pt);
  \end{tikzpicture}%
  \vspace{6pt}%
}

% ─── Clearance level badge ──────────────────────────────────────────
% Usage: \clearancebadge{4}
\newcommand{\clearancebadge}[1]{%
  \begin{tikzpicture}[baseline=(badge.base)]
    \node[
      draw=scpred,
      fill=scpred,
      line width=0.8pt,
      rounded corners=2pt,
      inner xsep=5pt,
      inner ysep=2pt,
      font=\sffamily\bfseries\scriptsize\color{white}
    ] (badge) {LEVEL #1};
  \end{tikzpicture}%
}

\endinput


% ─── Classification stamp (inline) ──────────────────────────────────
% Usage: \classificationstamp{TEXT}
\newcommand{\classificationstamp}[1]{%
  \begin{tikzpicture}[baseline=(text.base)]
    \node[
      draw=scpstamp,
      line width=1.5pt,
      rounded corners=1pt,
      inner sep=3pt,
      font=\sffamily\bfseries\scriptsize\color{scpstamp},
      opacity=0.85
    ] (text) {\MakeUppercase{#1}};
  \end{tikzpicture}%
}

% ─── Rotated classification stamp (for overlays) ────────────────────
% Usage: \rotatedstamp{TEXT}{angle}
\newcommand{\rotatedstamp}[2]{%
  \begin{tikzpicture}
    \node[
      draw=scpstamp,
      line width=2.5pt,
      rounded corners=2pt,
      rotate=#2,
      inner sep=6pt,
      font=\fontsize{18}{20}\selectfont\ttfamily\bfseries\color{scpstamp},
      opacity=0.7
    ] {\MakeUppercase{#1}};
  \end{tikzpicture}%
}

% ─── Document control number ────────────────────────────────────────
% Typewriter-style control number for headers or metadata
% Usage: \controlnumber{DOC-SCP-2026-0001}
\newcommand{\controlnumber}[1]{%
  {\ttfamily\scriptsize\color{scpgray}#1}%
}

% ─── Redaction bar ──────────────────────────────────────────────────
% Black bar for redacted text of arbitrary length
% Usage: \redacted{approximate width, e.g. 3cm}
\newcommand{\redacted}[1]{%
  \tikz[baseline=(bar.base)]{%
    \node[
      fill=scpblack,
      minimum width=#1,
      minimum height=1.1em,
      inner sep=0pt
    ] (bar) {};
  }%
}

% ─── Section divider ────────────────────────────────────────────────
% Red-accented horizontal rule
% Usage: \scpdivider
\newcommand{\scpdivider}{%
  \vspace{6pt}%
  \noindent
  \begin{tikzpicture}
    \draw[scpred, line width=0.8pt] (0,0) -- (\textwidth,0);
    \fill[scpred] (0.5\textwidth-3pt, -2pt) rectangle (0.5\textwidth+3pt, 2pt);
  \end{tikzpicture}%
  \vspace{6pt}%
}

% ─── Clearance level badge ──────────────────────────────────────────
% Usage: \clearancebadge{4}
\newcommand{\clearancebadge}[1]{%
  \begin{tikzpicture}[baseline=(badge.base)]
    \node[
      draw=scpred,
      fill=scpred,
      line width=0.8pt,
      rounded corners=2pt,
      inner xsep=5pt,
      inner ysep=2pt,
      font=\sffamily\bfseries\scriptsize\color{white}
    ] (badge) {LEVEL #1};
  \end{tikzpicture}%
}

\endinput
